\documentclass[10pt,a4paper]{amsart}
\usepackage{amsmath,amssymb,amsthm,graphicx}
% Retained minimal preamble for the Stacks Project chapter more-algebra.tex (More on Algebra).
% The chapter file pinned-original-source.tex carries only \input{preamble} on its first line, so no
% preamble text lives inside the chapter itself.  The Stacks Project uses one single shared preamble for
% every chapter; the snapshot stages it identically as sources/stacks-chow/preamble.tex and
% sources/stacks-etale-cohomology/preamble.tex.  The block below is copied mechanically, line for line,
% from sources/stacks-chow/preamble.tex lines 14-18 (the Xy-pic setup) and 154-217 (the native theorem
% declarations together with the mathematical macros).  Nothing was added, removed or redefined; the
% wrapper class and the amsmath/amssymb/amsthm/graphicx base packages are supplied by the assembler.
% Chapter: more-algebra.tex -- More on Algebra.
% For commutative diagrams we use Xy-pic
\usepackage[all]{xy}

% We use 2cell for 2-commutative diagrams.
\xyoption{2cell}

% Theorem environments.
%
\theoremstyle{plain}
\newtheorem{theorem}[subsection]{Theorem}
\newtheorem{proposition}[subsection]{Proposition}
\newtheorem{lemma}[subsection]{Lemma}

\theoremstyle{definition}
\newtheorem{definition}[subsection]{Definition}
\newtheorem{example}[subsection]{Example}
\newtheorem{exercise}[subsection]{Exercise}
\newtheorem{situation}[subsection]{Situation}

\theoremstyle{remark}
\newtheorem{remark}[subsection]{Remark}
\newtheorem{remarks}[subsection]{Remarks}

\numberwithin{equation}{subsection}

% Macros
%
\def\lim{\mathop{\mathrm{lim}}\nolimits}
\def\colim{\mathop{\mathrm{colim}}\nolimits}
\def\Spec{\mathop{\mathrm{Spec}}}
\def\Hom{\mathop{\mathrm{Hom}}\nolimits}
\def\Ext{\mathop{\mathrm{Ext}}\nolimits}
\def\SheafHom{\mathop{\mathcal{H}\!\mathit{om}}\nolimits}
\def\SheafExt{\mathop{\mathcal{E}\!\mathit{xt}}\nolimits}
\def\Sch{\mathit{Sch}}
\def\Mor{\mathop{\mathrm{Mor}}\nolimits}
\def\Ob{\mathop{\mathrm{Ob}}\nolimits}
\def\Sh{\mathop{\mathit{Sh}}\nolimits}
\def\NL{\mathop{N\!L}\nolimits}
\def\CH{\mathop{\mathrm{CH}}\nolimits}
\def\proetale{{pro\text{-}\acute{e}tale}}
\def\etale{{\acute{e}tale}}
\def\QCoh{\mathit{QCoh}}
\def\Ker{\mathop{\mathrm{Ker}}}
\def\Im{\mathop{\mathrm{Im}}}
\def\Coker{\mathop{\mathrm{Coker}}}
\def\Coim{\mathop{\mathrm{Coim}}}

% Boxtimes
%
\DeclareMathSymbol{\boxtimes}{\mathbin}{AMSa}{"02}

%
% Macros for moduli stacks/spaces
%
\def\QCohstack{\mathcal{QC}\!\mathit{oh}}
\def\Cohstack{\mathcal{C}\!\mathit{oh}}
\def\Spacesstack{\mathcal{S}\!\mathit{paces}}
\def\Quotfunctor{\mathrm{Quot}}
\def\Hilbfunctor{\mathrm{Hilb}}
\def\Curvesstack{\mathcal{C}\!\mathit{urves}}
\def\Polarizedstack{\mathcal{P}\!\mathit{olarized}}
\def\Complexesstack{\mathcal{C}\!\mathit{omplexes}}
% \Pic is the operator that assigns to X its picard group, usage \Pic(X)
% \Picardstack_{X/B} denotes the Picard stack of X over B
% \Picardfunctor_{X/B} denotes the Picard functor of X over B
\def\Pic{\mathop{\mathrm{Pic}}\nolimits}
\def\Picardstack{\mathcal{P}\!\mathit{ic}}
\def\Picardfunctor{\mathrm{Pic}}
\def\Deformationcategory{\mathcal{D}\!\mathit{ef}}

% --- appended 2026-10-12 (Chief): the project own definitions of the reference/slogan/history
% display environments, copied mechanically from the shared Stacks preamble lines 8-12
% (\usepackage{verbatim} plus the three \newenvironment lines, which the chapter preamble keeps
% verbatim from the project preamble). Nothing added, removed or redefined: these environments
% carry no mathematical content and the project itself typesets them as comments. Without these
% three lines a counted statement whose own text uses them fails with Environment slogan undefined.
\usepackage{verbatim}
\newenvironment{reference}{\comment}{\endcomment}
\newenvironment{slogan}{\comment}{\endcomment}
\newenvironment{history}{\comment}{\endcomment}
\usepackage{cleveref}
\crefname{originallemma}{original lemma}{original lemmas}
\Crefname{originallemma}{Original lemma}{Original lemmas}
\makeatletter
\newlabel{more-algebra-lemma-helper-G-ring}{{07PR}{source}{}{sourcepointer-1}{}}
\newlabel{more-algebra-lemma-helper-G-ring@cref}{{[originallemma][0][]07PR}{[1][0][]source}{}{}{}}
\makeatother
\crefname{originallemma}{original lemma}{original lemmas}
\Crefname{originallemma}{Original lemma}{Original lemmas}
\makeatletter
\newlabel{lemma-completion-regular}{{07NY}{source}{}{sourcepointer-2}{}}
\newlabel{lemma-completion-regular@cref}{{[originallemma][0][]07NY}{[1][0][]source}{}{}{}}
\makeatother
\crefname{originallemma}{original lemma}{original lemmas}
\Crefname{originallemma}{Original lemma}{Original lemmas}
\makeatletter
\newlabel{algebra-lemma-completion-finite-extension}{{07N9}{source}{}{sourcepointer-3}{}}
\newlabel{algebra-lemma-completion-finite-extension@cref}{{[originallemma][0][]07N9}{[1][0][]source}{}{}{}}
\makeatother
\crefname{originallemma}{original lemma}{original lemmas}
\Crefname{originallemma}{Original lemma}{Original lemmas}
\makeatletter
\newlabel{lemma-find-D}{{07PH}{source}{}{sourcepointer-4}{}}
\newlabel{lemma-find-D@cref}{{[originallemma][0][]07PH}{[1][0][]source}{}{}{}}
\makeatother
\crefname{originallemma}{original lemma}{original lemmas}
\Crefname{originallemma}{Original lemma}{Original lemmas}
\makeatletter
\newlabel{lemma-derivation-extends}{{07PE}{source}{}{sourcepointer-5}{}}
\newlabel{lemma-derivation-extends@cref}{{[originallemma][0][]07PE}{[1][0][]source}{}{}{}}
\makeatother
\crefname{originallemma}{original lemma}{original lemmas}
\Crefname{originallemma}{Original lemma}{Original lemmas}
\makeatletter
\newlabel{lemma-degree-p-extension-regular}{{07PG}{source}{}{sourcepointer-6}{}}
\newlabel{lemma-degree-p-extension-regular@cref}{{[originallemma][0][]07PG}{[1][0][]source}{}{}{}}
\makeatother
\title{Geometric regularity of a completed localisation of a power series ring over a polynomial ring}
\author{The Stacks Project Authors}
\date{stacks-project:more-algebra.tex:local-snapshot-20261012; GNU FDL 1.2 or later}
\begin{document}
\maketitle
\setcounter{subsection}{1}
\noindent\textit{Editorial note.} All delivered mathematical text is the author's own, copied from the pinned Stacks Project chapter file more-algebra.tex; no statement or proof step was rewritten, reordered or removed.\par\medskip
\noindent\textit{Editorial note.} The chapter file pinned-original-source.tex carries only \textbackslash{}input\{preamble\} on its first line, so no preamble text lives inside the chapter. The retained file sources/stacks-more-algebra/retained-preamble.tex is a mechanical, line-for-line copy of the project's single shared preamble as staged in the snapshot (Xy-pic setup plus theorem declarations and mathematical macros). Nothing was added, removed or redefined, and the excerpt serves only as preamble, never as a span coordinate.\par\medskip
\noindent\textit{Editorial note.} Both counted spans start at the \textbackslash{}begin line of their environment, so the statement and the author's proof are delivered whole; surrounding prose of the section is deliberately outside the counted spans and is carried as standard prerequisites instead.\par\medskip
\noindent\textit{Editorial note.} The citation fields (source\_id, source\_version, authors, source\_url, license.name, source\_title) are identical to those of every other card this lane has produced from this chapter, so the batch selector treats them as one source.\par\medskip
\noindent\textit{Allowed standard prerequisites (editorial).}\begin{itemize}
\item The completion of a regular local ring is regular. Lemma \textbackslash{}ref\{lemma-completion-regular\}, cited at proof line 12317. Gives the base case of the induction.
\item Completion of a finite extension of Noetherian complete local rings is a product of completions. Algebra, Lemma \textbackslash{}ref\{algebra-lemma-completion-finite-extension\}, cited at proof lines 12332-12333. Identifies the base change with the completed localisation of the finite algebra realising the field extension.
\item Existence of a derivation of a finite algebra that does not vanish on a prescribed element. Lemma \textbackslash{}ref\{lemma-find-D\}, cited at proof line 12341. Produces a derivation D of B with D(f) nonzero, so that D(f) becomes a unit after inverting denominators.
\item A derivation extends to a localisation and to a completion. Lemma \textbackslash{}ref\{lemma-derivation-extends\}, cited at proof lines 12343-12346. Extends D to the completed localisation tensored with the intermediate field, where D(f) is a unit.
\item A degree-p extension obtained by adjoining a p-th root is regular when a derivation sends that root's argument to a unit. Lemma \textbackslash{}ref\{lemma-degree-p-extension-regular\}, cited at proof lines 12350-12351. Concludes regularity of the base change after adjoining the p-th root.
\end{itemize}\medskip
\ifdefined\hypertarget\hypertarget{sourcepointer-1}{}\fi
\noindent\textit{Source reference.} Original lemma 07PR: focused-seeds/AGP/sources/stacks-snapshot-20261011/tags.txt, source lines 9806--9806. The lemma's own canonical tag, quoted so that readers can resolve the label inside the delivered text; it points at the counted statement itself and carries no proof content.\par\medskip
\ifdefined\hypertarget\hypertarget{sourcepointer-2}{}\fi
\noindent\textit{Source reference.} Original lemma 07NY: focused-seeds/AGP/sources/stacks-snapshot-20261011/tags.txt, source lines 9778--9778. Navigation pointer to a statement the counted proof refers to but does not contain; nothing load-bearing is delegated to it.\par\medskip
\ifdefined\hypertarget\hypertarget{sourcepointer-3}{}\fi
\noindent\textit{Source reference.} Original lemma 07N9: focused-seeds/AGP/sources/stacks-snapshot-20261011/tags.txt, source lines 9754--9754. Navigation pointer to a statement the counted proof refers to but does not contain; nothing load-bearing is delegated to it.\par\medskip
\ifdefined\hypertarget\hypertarget{sourcepointer-4}{}\fi
\noindent\textit{Source reference.} Original lemma 07PH: focused-seeds/AGP/sources/stacks-snapshot-20261011/tags.txt, source lines 9797--9797. Navigation pointer to a statement the counted proof refers to but does not contain; nothing load-bearing is delegated to it.\par\medskip
\ifdefined\hypertarget\hypertarget{sourcepointer-5}{}\fi
\noindent\textit{Source reference.} Original lemma 07PE: focused-seeds/AGP/sources/stacks-snapshot-20261011/tags.txt, source lines 9794--9794. Navigation pointer to a statement the counted proof refers to but does not contain; nothing load-bearing is delegated to it.\par\medskip
\ifdefined\hypertarget\hypertarget{sourcepointer-6}{}\fi
\noindent\textit{Source reference.} Original lemma 07PG: focused-seeds/AGP/sources/stacks-snapshot-20261011/tags.txt, source lines 9796--9796. Navigation pointer to a statement the counted proof refers to but does not contain; nothing load-bearing is delegated to it.\par\medskip
\begin{lemma}
\label{lemma-helper-G-ring}
Let $k$ be a field of characteristic $p$.
Let $A = k[[x_1, \ldots, x_n]][y_1, \ldots, y_m]$ and denote $K$
the fraction field of $A$.
Let $\mathfrak p \subset A$ be a prime. Then
$A_\mathfrak p^\wedge \otimes_A K$ is geometrically regular over $K$.
\end{lemma}

\begin{proof}
Let $L/K$ be a finite purely inseparable field extension.
We will show by induction on $[L : K]$ that $A_\mathfrak p^\wedge \otimes L$
is regular. The base case is $L = K$: as $A$ is regular,
$A_\mathfrak p^\wedge$ is regular (Lemma \ref{lemma-completion-regular}),
hence the localization $A_\mathfrak p^\wedge \otimes K$ is regular.
Let $K \subset M \subset L$ be a subfield such that
$L$ is a degree $p$ extension of $M$ obtained by adjoining a $p$th root
of an element $f \in M$. Let $B$ be a finite $A$-subalgebra
of $M$ with fraction field $M$. Clearing denominators, we may and do assume
$f \in B$. Set $C = B[z]/(z^p -f)$ and note that $B \subset C$
is finite and that the fraction field of $C$ is $L$. Since
$A \subset B \subset C$ are finite and $L/M/K$ are purely
inseparable we see that for every element of $B$ or $C$ some power of
it lies in $A$. Hence there is a unique prime $\mathfrak r \subset B$,
resp.\ $\mathfrak q \subset C$ lying over $\mathfrak p$. Note that
$$
A_\mathfrak p^\wedge \otimes_A M = B_\mathfrak r^\wedge \otimes_B M
$$
see Algebra, Lemma \ref{algebra-lemma-completion-finite-extension}.
By induction we know that this ring is regular. In the same manner we have
$$
A_\mathfrak p^\wedge \otimes_A L =
C_\mathfrak r^\wedge \otimes_C L =
B_\mathfrak r^\wedge \otimes_B M[z]/(z^p - f)
$$
the last equality because the completion of
$C = B[z]/(z^p - f)$ equals $B_\mathfrak r^\wedge[z]/(z^p -f)$.
By Lemma \ref{lemma-find-D} we know there exists a derivation
$D : B \to B$ such that $D(f) \not = 0$. In other words, $g = D(f)$
is a unit in $M$! By Lemma \ref{lemma-derivation-extends}
$D$ extends to a derivation of $B_\mathfrak r$, $B_\mathfrak r^\wedge$
and $B_\mathfrak r^\wedge \otimes_B M$ (successively extending through a
localization, a completion, and a localization). Since it is an
extension we end up with a derivation of $B_\mathfrak r^\wedge \otimes_B M$
which maps $f$ to $g$ and $g$ is a unit of the ring
$B_\mathfrak r^\wedge \otimes_B M$.
Hence $A_\mathfrak p^\wedge \otimes_A L$ is regular by
Lemma \ref{lemma-degree-p-extension-regular} and we win.
\end{proof}

\end{document}
